\begin{abstract}
Public leaderboards rank models by aggregating scores on several sub-datasets, and a developer can try to raise a chosen model's rank by training on data related to some of them. How much such targeted investment is needed to reach first place, and how does it depend on the aggregation rule? We measure investment as the score points added on chosen sub-datasets up to a cap, by the number of sub-datasets touched ($K$), the total points ($L_1$) and the largest single jump ($L_\infty$), and compute the minimum exactly for ten rules (mean, power and trimmed means, median, minimum, Borda, plurality, Condorcet winner), with closed forms for the mean, minimum and median. On the Open LLM Leaderboard (4496 models, 6 tasks) and HELM MMLU v1.13.0 (79 models, 57 subjects) we find that (i) with no cap beyond the scale maximum, one sub-dataset with enough headroom suffices under the mean for most targets up to rank 500 on OLLB, but the points needed are $n$ times the gap in mean score; (ii) in points the averaging rules are the most expensive and the median, minimum and plurality rules far cheaper, while ordinal majority rules need the most tasks, so no rule is cheapest on every measure and the order of rules changes with the dataset and the cap; (iii) targeting saves points over a uniform improvement mainly for rules other than the mean (cost 3--51\% of the uniform cost for median, minimum and plurality, against 93--100\% for the mean at the scale maximum); (iv) the answer depends on the assumed cap: under a cap taken from observed fine-tune gains on one leaderboard, 7.5\% of OLLB targets can become first under the mean, and 1.4--12\% for the mean across the cap definitions we tried. Investment is measured in score points, not in data or compute, and the conclusions are conditional on the cap, a single deviating developer and the two leaderboards; ordinal rules are solved on a stratified sample.
\end{abstract}
